\chapter{Membagi Rata dan Pembagian}\label{ch:g3:division}

Dua puluh kelereng yang dibagi rata kepada empat anak: berapa untuk
masing-masing? Pembagian menjawab semua pertanyaan tentang membagi
rata --- dan pertanyaan tentang mengelompokkan juga --- dan ia
kembaran perkalian: yang satu membatalkan yang lain.

\section{Dua wajah pembagian}

\begin{definition}[Pembagian]\label{def:g3:division:def}
Pembagian $20 \div 4$ menjawab dua macam pertanyaan:
\begin{itemize}
  \item \emph{membagi rata}: $20$ kelereng dibagi rata kepada $4$
        anak --- berapa kelereng untuk \emph{masing-masing}? ($5$
        untuk tiap anak);
  \item \emph{mengelompokkan}: $20$ kelereng dimasukkan ke kantong
        berisi $4$ --- berapa \emph{kantong}? ($5$ kantong).
\end{itemize}
Kedua pertanyaan itu punya jawaban yang sama, karena keduanya
membatalkan perkalian yang sama: $4 \times 5 = 20$.
\end{definition}

\begin{omfigure}
\begin{tikzpicture}[scale=0.55]
  \foreach \g in {0,1,2,3} {
    \draw[omExo, rounded corners] (\g*4.9-0.65,-0.7) rectangle
      (\g*4.9+3.85,0.7);
    \foreach \c in {0,...,4} \fill[omDef] (\g*4.9+\c*0.8,0) circle (0.27);
  }
  \node[below, font=\small] at (8.95,-1.15)
    {$20$ kelereng, $4$ kelompok sama besar: $5$ di tiap kelompok};
\end{tikzpicture}

{\small Membagi $20$ menjadi $4$ kelompok yang sama besar.
Mengelompokkan ke dalam kantong berisi $4$ malah memberi $5$
kelompok --- pembagian yang sama juga.}
\end{omfigure}

\begin{method}[Membagi dengan tabel perkalian]\label{met:g3:division:tables}
Untuk menghitung $42 \div 6$, tanyakanlah: ``$6$ dikali \emph{berapa}
sama dengan $42$?'' Tabel $6$ menjawabnya: $6 \times 7 = 42$, jadi
$42 \div 6 = 7$. Setiap pertanyaan pembagian adalah pertanyaan
perkalian yang dibaca terbalik.
\end{method}

\section{Ketika pembagiannya tidak pas}

\begin{definition}[Sisa]\label{def:g3:division:remainder}
Membagi rata $23$ kelereng kepada $4$ anak memberi $5$ kelereng
untuk masing-masing --- dan $3$ kelereng \emph{tersisa}, tidak cukup
untuk satu putaran lagi. Kita menulis
\[
23 = 4 \times 5 + 3 ,
\]
lalu berkata: \emph{hasil bagi}\index{hasil bagi} adalah $5$, dan
\emph{sisa}\index{sisa} adalah $3$. Sisa selalu lebih kecil daripada
banyaknya anak --- kalau tidak, satu putaran pembagian lagi masih
mungkin dilakukan.
\end{definition}

\begin{omfigure}
\begin{tikzpicture}[scale=0.55]
  \foreach \g in {0,1,2,3} {
    \draw[omExo, rounded corners] (\g*4.9-0.65,-0.7) rectangle
      (\g*4.9+3.85,0.7);
    \foreach \c in {0,...,4} \fill[omDef] (\g*4.9+\c*0.8,0) circle (0.27);
  }
  \foreach \c in {0,1,2} \fill[omThm] (19.7+\c*0.8,0) circle (0.27);
  \node[above, font=\footnotesize, omThm] at (20.5,0.5) {tersisa};
  \node[below, font=\small] at (10.2,-1.15)
    {$23 = 4 \times 5 + 3$: hasil bagi $5$, sisa $3$ (berwarna merah)};
\end{tikzpicture}

{\small Membagi $23$ kepada $4$: setelah $5$ putaran, $3$ kelereng
tersisa --- terlalu sedikit untuk putaran ke-$6$.}
\end{omfigure}

\begin{example}\label{ex:g3:division:remainder}
Bagilah $58$ dengan $7$. Telusuri tabel $7$:
$7 \times 8 = 56$ adalah \omterm{def:g2:mult:def}{hasil kali} terbesar yang muat di bawah $58$
($7 \times 9 = 63$ sudah terlalu besar). Jadi
\[
58 = 7 \times 8 + 2 :
\]
\omterm{def:g3:division:remainder}{hasil bagi} $8$, \omterm{def:g3:division:remainder}{sisa} $2$. Periksa: $56 + 2 = 58$, dan $2 < 7$.
\end{example}

\begin{method}[Menafsirkan sisa]\label{met:g3:division:interpret}
Setelah membagi, kembalilah ke ceritanya:
\begin{enumerate}
  \item ``berapa untuk masing-masing / berapa kelompok yang penuh?''
        --- \omterm{def:g3:division:remainder}{hasil baginya};
  \item ``berapa yang tersisa?'' --- \omterm{def:g3:division:remainder}{sisanya};
  \item ``berapa kotak untuk memuat semuanya?'' --- \omterm{def:g3:division:remainder}{hasil baginya}
        \emph{ditambah satu} kalau \omterm{def:g3:division:remainder}{sisanya} bukan \omterm{def:g1:counting:zero}{nol}.
\end{enumerate}
\end{method}

\begin{example}\label{ex:g3:division:interpret}
$58$ murid pergi berkano, $7$ orang tiap kano. $58 = 7 \times 8 + 2$:
delapan kano penuh, dan $2$ murid tersisa --- mereka juga perlu
kano. Jadi $9$ kano harus dipesan.
\end{example}

\section{Setengah, sepertiga, seperempat}

\begin{definition}[Setengah, sepertiga, seperempat]\label{def:g3:division:half}
\emph{Setengah}\index{setengah} dari sebuah bilangan adalah hasil
membaginya dengan $2$; \emph{sepertiga}, dengan $3$;
\emph{seperempat}\index{seperempat}, dengan $4$. Misalnya setengah
dari $18$ adalah $9$, sepertiga dari $18$ adalah $6$, dan seperempat
dari $36$ adalah $9$.
\end{definition}

\begin{example}[Kembar dan setengah selalu berpasangan]\label{ex:g3:division:double}
Setengah dari $46$: karena $46 = 40 + 6$, setengah dari tiap bagian
memberi $20 + 3 = 23$. Periksa dengan kembarnya: $23 + 23 = 46$.
Hafal kembar di luar kepala ($25 \to 50$, $45 \to 90$,
$150 \to 300$) membuat setengah menjadi seketika.
\end{example}

\section{Latihan}

\begin{exercise}[$\star$]\label{exo:g3:division:1}
Hitunglah dengan membaca tabel secara terbalik: $36 \div 4$; \quad
$63 \div 9$; \quad $40 \div 5$; \quad $56 \div 8$.
\end{exercise}

\begin{exercise}[$\star$]\label{exo:g3:division:2}
Untuk setiap pembagian, tulislah kesamaan $a = b \times q + r$ lalu
sebutkan \omterm{def:g3:division:remainder}{hasil bagi} dan sisanya: $17 \div 5$; \quad $30 \div 4$; \quad
$50 \div 6$.
\end{exercise}

\begin{exercise}[$\star$]\label{exo:g3:division:3}
$28$ kelereng dibagi rata kepada $6$ anak. Berapa kelereng untuk
masing-masing, dan berapa yang tersisa? Gambarlah pembagiannya kalau membantu.
\end{exercise}

\begin{exercise}[$\star$]\label{exo:g3:division:4}
Hitunglah: setengah dari $68$; sepertiga dari $27$; \omterm{def:g3:division:half}{seperempat} dari
$48$; setengah dari $150$.
\end{exercise}

\begin{exercise}[$\star$]\label{exo:g3:division:5}
Benar atau salah? Periksalah dengan perkalian.
``$45 \div 5 = 9$''; \quad ``$52 \div 6 = 8$ \omterm{def:g3:division:remainder}{sisa} $4$'';
\quad ``$70 \div 8 = 9$ \omterm{def:g3:division:remainder}{sisa} $2$''.
\end{exercise}

\begin{exercise}[$\star$]\label{exo:g3:division:6}
$32$ foto ditempelkan ke dalam album, $4$ foto per halaman. Berapa
halaman yang terpakai? Ini wajah pembagian yang mana (membagi rata atau mengelompokkan)?
\end{exercise}

\begin{exercise}[$\star$]\label{exo:g3:division:7}
Seutas tali yang panjangnya $75$ cm dipotong menjadi potongan $9$ cm.
Berapa potongan yang utuh, dan berapa panjang sisa potongannya?
\end{exercise}

\begin{exercise}[$\star$]\label{exo:g3:division:8}
$50$ anak pergi bertamasya dengan mobil berisi $4$ tempat duduk.
Berapa mobil yang diperlukan? (Hati-hati: semua harus terangkut!)
\end{exercise}

\begin{exercise}[$\star$]\label{exo:g3:division:9}
Carilah bilangan yang hilang: $? \div 7 = 6$; \quad
$54 \div ? = 6$; \quad setengah dari $?$ adalah $35$.
\end{exercise}

\begin{exercise}[$\star\star$]\label{exo:g3:division:10}
Amina membagi $47$ stikernya kepada $5$ temannya, serata mungkin,
lalu menyimpan sisanya untuk dirinya sendiri. Berapa stiker yang
didapat setiap teman, dan berapa yang disimpan Amina?
\end{exercise}

\begin{exercise}[$\star\star$]\label{exo:g3:division:11}
Ketika sebuah bilangan dibagi $6$, \omterm{def:g3:division:remainder}{sisa} apa saja yang mungkin?
Berapa bilangan terbesar yang pembagiannya dengan $6$ menghasilkan
\omterm{def:g3:division:remainder}{hasil bagi} $7$?

\end{exercise}
